%======================================================================
\chapter{Theorem 15: ladders}\label{ch:ladders}
%======================================================================

It remains to prove Theorem~\ref{thm:ho15}: inside any neighbourhood $U$ of a separator $M$ of the bottom of $G^\star$ from the top there is a set with a
stairwell structure of odd height. Hoehn and Oversteegen first replace $M$ by a nicer separator $M'$ (their Claim~15.1, ``left to the reader''), make it irreducible,
and then build the stairwell over the arcs of a subdivision, with a zig-zag added wherever two curves of $M'$ are born or die. This chapter replaces Claim~15.1
and the irreducibility by an explicit combinatorial object, a \emph{ladder}, read off from the signs of one function on a grid. The odd counts come from
sign changes, and the parity argument of Hoehn and Oversteegen for (S4) becomes a statement about parities of list lengths.

\section{Ladders}

Fix a decomposition $\mathcal S$ of the PL graph $G$ into oriented edges $s=(s_1,s_2)$, and write $s(\mu)=s_1+\mu(s_2-s_1)$ for $\mu\in\I$.

\begin{definition}[Ladder]\label{def:ladder}
A \emph{ladder} on $G$ consists of a decomposition $\mathcal S$ of $G$, a list $y(v)=(y_1(v),\dots,y_{j(v)}(v))$ of heights in $(0,1)$ of \emph{odd} length $j(v)$ at each vertex $v$, and
for each edge $s$ an \emph{event}, which is either
\begin{itemize}
\item \emph{none}: then $j(s_1)=j(s_2)$; or
\item a \emph{wedge} $(m,\lambda,h)$ at $s_2$: then $j(s_2)=j(s_1)+2$, $1\le m\le j(s_1)+1$, $\lambda,h\in(0,1)$. The wedge occupies places $m$ and $m+1$ at $s_2$, and its tip is
$(s(\lambda),h)$.
\end{itemize}
The \emph{height} of the ladder is the largest $j(v)$.
\end{definition}

\formal{\lean{HPC.HO16.IsLadder}, \lean{HPC.HO16.ladderY}, \lean{HPC.HO16.ladderHeight} (in \lean{HPC/HO16/LadderDefs.lean}).}

\emph{The sheets.} Let $s$ be an edge and $j=j(s_1)$.
\begin{itemize}
\item No event: for $1\le i\le j$, sheet $i$ contains the segment from $(s_1,y_i(s_1))$ to $(s_2,y_i(s_2))$.
\item A wedge $(m,\lambda,h)$: for $i<m$ as before. For $m\le i\le j$, let $L_i$ be the line from $(s_1,y_i(s_1))$ to $(s_2,y_{i+2}(s_2))$, parametrized by $\mu\in\I$, and put
\[
l_i=\lambda\,\frac{2(j-i)+1}{2(j-m)+3},\qquad r_i=\lambda\,\frac{2(j-i)+2}{2(j-m)+3},
\]
so that $0<l_j<r_j<l_{j-1}<\dots<l_m<r_m<\lambda$. The \emph{zig-zag} of curve $i$ consists of $C_i^1=L_i[0,r_i]$ in sheet $i$, $C_i^2=L_i[l_i,r_i]$ in sheet $i+1$, and $C_i^3=L_i[l_i,1]$ in
sheet $i+2$. Finally the two halves of the wedge, the segments from the tip to $(s_2,y_m(s_2))$ and to $(s_2,y_{m+1}(s_2))$, go into sheets $m$ and $m+1$.
\end{itemize}
Sheet $i$ is the union of its pieces over all edges.

\formal{\lean{HPC.HO16.ladderPiece}, \lean{HPC.HO16.ladderSheet}; the zig-zag parameters are \lean{HPC.HO16.zzL}, \lean{HPC.HO16.zzR} and the lines
\lean{HPC.HO16.zzLine}.}

Over a wedge edge, a curve $i\ge m$ entering at rank $i$ leaves at rank $i+2$: it climbs two sheets by running forward in sheet $i$, back in sheet $i+1$ and forward in
sheet $i+2$. The three pieces lie on one line and overlap. Definitions~\ref{def:straight}--\ref{def:stairwell} do not forbid sheets to meet, and the end sets
come out of the construction, so nothing needs to be separated. Compared with Hoehn and Oversteegen's picture, the wedge is put at the end of the edge and the zig-zags
before its tip.

\begin{example}\label{ex:ladder}
$G=[0,2]\subseteq\R$ with vertices $0,1,2$ and lists $(\frac12)$, $(\frac12)$, $(\frac14,\frac12,\frac34)$; no event on $(0,1)$ and a wedge $(1,\frac12,\frac38)$ at $2$ on $(1,2)$. Here $j=1$, $m=1$,
$l_1=\frac16$, $r_1=\frac13$ in the parameter of $[1,2]$. Then $\pi(S_1)=[0,\frac43]\cup[\frac32,2]$, $\pi(S_2)=[\frac76,\frac43]\cup[\frac32,2]$, $\pi(S_3)=[\frac76,2]$, $\pi(\alpha_2)=\set{\frac43,\frac32}$,
$\pi(\alpha_3)=\set{\frac76}$: a stairwell structure of height~$3$.
\end{example}

\formal{\lean{HPC.HO16.ladder_example}, \lean{HPC.HO16.ladder_example_height}, \lean{HPC.HO16.ladder_example_stairwell}.}

\section{A ladder is a stairwell}

\begin{theorem}\label{thm:ladder-stairwell}
Let a ladder on $G$ with at least one edge have height $k$. Then $k$ is odd, and the sheets $S_1,\dots,S_k$ form a stairwell structure of height $k$ for their union.
\end{theorem}

\begin{proof}
$k$ is the length of some list, hence odd.

\emph{The end sets.} On a wedge edge write $L_i(\mu)$ for the point of $L_i$ at parameter $\mu$. Let $\alpha_i$ consist of the points $L_{i-1}(r_{i-1})$ (the common right end of $C^1_{i-1}$ and
$C^2_{i-1}$, if $m\le i-1\le j$), $L_{i-2}(l_{i-2})$ (the common left end of $C^2_{i-2}$ and $C^3_{i-2}$, if $m\le i-2\le j$), and the tip (if $i=m+1$), over all wedge edges. Then
$\alpha_1=\emptyset$, and $\alpha_{k+1}=\emptyset$ because a wedge edge has $j(s_1)\le k-2$.

(S1) Over an edge, the pieces of sheet $i$ lie over the parameter intervals $[0,r_i]$, $[l_{i-1},r_{i-1}]$, $[l_{i-2},1]$ and, for $i\in\set{m,m+1}$, $[\lambda,1]$, which are pairwise disjoint;
over a vertex $v$ every edge at $v$ contributes the same point $(v,y_i(v))$ if $i\le j(v)$, and none otherwise. So $\pi$ is injective on each sheet, and $\pi(S_i)$ is a finite union of
nondegenerate segments, a PL graph (Proposition~\ref{prop:pl-library}(a)). A vertex with $j(v)=k$ gives every sheet a point.

(S2) A vertex $v$ is interior to $\pi(S_i)$ if $i\le j(v)$ (every edge at $v$ has a piece of sheet $i$ starting at $v$) and outside it otherwise. So $\rfr\pi(S_i)$ consists of the parameters
where pieces stop inside edges: $r_i$, $l_{i-1}$, $r_{i-1}$, $l_{i-2}$ and the tip. These are exactly the points of $\alpha_i\cup\alpha_{i+1}$, and they lie in $S_i$.

(S3) At each point of $\beta_i=\alpha_{i+1}$ the two sheets leave in the same direction: $C_i^1$ and $C_i^2$ both end at $r_i$ from the left; $C_{i-1}^2$ and $C_{i-1}^3$ both start at
$l_{i-1}$ to the right; the two wedge halves both start at the tip. Every other piece of the two sheets on that edge is at positive distance.

(S4) Let $C$ be a component of $G\setminus\pi(S_i)$. Inside one edge, $C$ is one of the gaps $(r_i,l_{i-1})$ (both ends in $\pi(\beta_i)$), $(r_{i-1},l_{i-2})$ (both ends in $\pi(\alpha_i)$), or a gap
next to the tip (both ends in $\pi(\beta_m)$, respectively in $\pi(\alpha_{m+1})$). Otherwise $C$ contains vertices, all with $j(v)<i$. Its frontier lies on wedge edges $s$ with $j(s_1)<i\le
j(s_1)+2$, at the first piece of sheet $i$ starting from $s_1$, and that end point is in $\beta_i$ if $i-j(s_1)=1$ and in $\alpha_i$ if $i-j(s_1)=2$. \emph{All $j(v)$ are odd}, so which case occurs
depends only on the parity of $i$: $\rfr C\subseteq\pi(\beta_i)$ for $i$ even and $\rfr C\subseteq\pi(\alpha_i)$ for $i$ odd. As $\pi(\alpha_i)\cap\pi(\beta_i)=\emptyset$, both halves of (S4) follow.

(S5) All end points lie over interior points of edges, which are ordinary, at pairwise different parameters.
\end{proof}

\formal{\lean{HPC.HO16.ho16_stairwell_of_ladder}, with $\alpha$ given by \lean{HPC.HO16.ladderAlphaSet}; the parity argument for (S4) is \lean{HPC.HO16.ladder_S4}.}

This is Hoehn and Oversteegen's parity argument of p.~198 (``we must have that $j_1$ and $j_n$ have the same parity and, hence, $j_1=j_n$''). Connectedness of $G$ is not
needed, the lists need not be sorted, and the heights need not be distinct.

\section{A ladder inside a neighbourhood of a separator}

\begin{lemma}[A sign function]\label{lem:sign-fn}
Let $G$ be a nonempty compact set, $M\subseteq G\times(0,1)$ separate $G\times\set0$ from $G\times\set1$ in $G^\star$, and $U\supseteq M$ open. There are a continuous $u\colon E\times\R\to\R$ and $\delta>0$
with $u\le0$ on $G\times\set0$, $u>0$ on $G\times\set1$, and: if $p\in G^\star$ is within $\delta$ of a point $c_1$ with $u(c_1)\le0$ and of a point $c_2$ with $u(c_2)>0$, then $p\in U\cap G\times(0,1)$.
\end{lemma}

\begin{proof}
Let $O_0\supseteq G\times\set0$ and $O_1\supseteq G\times\set1$ be the open sets of the separation. $K=G^\star\setminus(U\cap E\times(0,1))$ is compact and lies in $G^\star\setminus M\subseteq O_0\cup O_1$. Put
$K_0=K\setminus O_1\supseteq G\times\set0$ and $K_1=K\setminus O_0\supseteq G\times\set1$: disjoint compact sets with union $K$. Let $u=\dist(\cdot,K_0)-\dist(\cdot,K_1)$, a $2$-Lipschitz function, negative on
$K_0$ and positive on $K_1$, so $\omega=\min_K\abs u>0$. Note $G\times\set0\subseteq K_0$: it lies in $K$, in $O_0$, and hence not in $O_1$, as $O_0\cap O_1$ misses $G^\star\setminus M$; likewise
$G\times\set1\subseteq K_1$. Take $\delta=\omega/4$. If $p$ is as in the statement, then $\abs{u(p)-u(c_i)}<2\delta=\omega/2$. If $u(p)\ge0$, then $u(p)<u(c_1)+\omega/2\le\omega/2$; if $u(p)<0$,
then $u(p)>u(c_2)-\omega/2>-\omega/2$. So $\abs{u(p)}<\omega$, hence $p\notin K$, that is, $p\in U\cap G\times(0,1)$.
\end{proof}

\formal{\lean{HPC.HO16.exists_sign_fn}.}

\begin{theorem}[Replacing Claim~15.1]\label{thm:exists-ladder}
Let $G$ be a connected PL graph, $M\subseteq G\times(0,1)$ separate $G\times\set0$ from $G\times\set1$ in $G^\star$, and $U\supseteq M$ open. There is a ladder on $G$ with at least one edge all of
whose sheets lie in $U\cap G\times(0,1)$.
\end{theorem}

\begin{proof}
Take $u$ and $\delta$ from Lemma~\ref{lem:sign-fn}. Choose a decomposition $\mathcal S_0$ of $G$ with edges shorter than $\delta$, and $N$ with $1/N<\delta$; put $t_q=q/N$. At a vertex $v$ let
$\sigma(v,q)=-$ if $u(v,t_q)\le0$ and $+$ otherwise; so $\sigma(v,0)=-$ and $\sigma(v,N)=+$.

\emph{Heights at vertices.} Let $y(v)$ be the list of midpoints $\frac12(t_q+t_{q+1})$ of those $q$ with $\sigma(v,q)\neq\sigma(v,q+1)$, in increasing order. Since the signs go from $-$ to $+$, there is an
odd number of them.

\emph{One flip per edge.} For an edge $(v,w)$ of $\mathcal S_0$ let $q_1<\dots<q_r$ be the indices with $\sigma(v,q)\neq\sigma(w,q)$. Subdivide the edge into $r$ edges (one if $r=0$), and give the $p$-th new
vertex the sign sequence of $v$ with $q_1,\dots,q_p$ flipped (so the last one has that of $w$), and the list of midpoints of its sign changes. Each new edge flips one sign $q$, which changes the
sign changes at $q-1\mid q$ and at $q\mid q+1$:
\begin{itemize}
\item one appears and one disappears: one curve moves from $\frac12(t_{q-1}+t_q)$ to $\frac12(t_q+t_{q+1})$ or back, and no other curve moves; no event;
\item two appear: a wedge at the far end, with $\lambda=\frac12$, $h=t_q$ and $m$ the rank of the new lower curve;
\item two disappear: a wedge at the near end; orient the edge backwards.
\end{itemize}

\emph{The sheets lie in $U$.} At every new vertex the sign at each level $t_q$ is that of $v$ or of $w$. A curve at height $\frac12(t_q+t_{q+1})$ at a vertex means a sign change between $t_q$ and $t_{q+1}$
there. So every piece of a sheet lies in a set $s_0\times[t_q,t_{q+1}]$ ($s_0$ the original edge) whose corners carry both signs, or, near a flip at $q$, in $s_0\times[t_{q-1},t_{q+1}]$, where
$(v,t_q)$ and $(w,t_q)$ carry both signs. Every point of such a set is within $\max(\abs{s_0},1/N)<\delta$ of each of its corners (in the sup-norm on $E\times\R$), so by Lemma~\ref{lem:sign-fn} it lies in
$U\cap G\times(0,1)$.
\end{proof}

\formal{\lean{HPC.HO16.exists_ladder} (in \lean{HPC/HO16/LadderSep.lean}, \lean{HPC/HO16/LadderSubdiv.lean}, \lean{HPC/HO16/LadderSign.lean},
\lean{HPC/HO16/LadderExists.lean}). Connectedness of $G$ is used only for $G\neq\emptyset$.}

\section{The end of the proof}

\begin{proof}[Proof of Theorem~\ref{thm:ho15}]
Theorem~\ref{thm:exists-ladder} gives a ladder whose sheets lie in $U\cap G\times(0,1)$, and Theorem~\ref{thm:ladder-stairwell} makes their union a set with a stairwell structure of odd height.
\end{proof}

\formal{\lean{HPC.HO16.ho16_theorem15}.}

\begin{corollary}\label{cor:final}
Theorem~\ref{thm:ho20} holds; hence Theorem~\ref{thm:ho1} holds without hypothesis, and every nondegenerate homogeneous indecomposable plane continuum is homeomorphic to the
pseudo-arc (Theorem~\ref{thm:indec-pseudoarc}).
\end{corollary}

\begin{proof}
Theorem~\ref{thm:20-from-15} and Theorem~\ref{thm:ho15}.
\end{proof}

\formal{\lean{HPC.HO16.ho16_theorem20}, \lean{HPC.HO16.ho16_indec_homogeneous_pseudoArc}, \lean{HPC.HO16.ho16_indec_homogeneous_iff}.}

\section{On Hoehn and Oversteegen's proof}

For the reader comparing with the paper (pp.~194--198):
\begin{itemize}
\item Claim~15.1, the existence of a separator $M'$ with properties (1), (3), (5), (6), is left to the reader. The passage to an irreducible separator keeps (1) and (3), but this is not argued.
In the definition of a side wedge, ``$C_2$ are arcs which both have $x$ as an endpoint'' should read $(x,y)$. And in the paragraph on property~(4), the reason given for ``$(x,y)$ cannot be an
endpoint of $M'$'' (property~(3)) does not apply at points of $T$; irreducibility gives the step.
\item The odd number of points of $M'$ over each $a\in Z'$ (p.~196) is not argued. Here it comes from sign changes.
\item In the zig-zag construction nothing excludes $b\in J_j$; the arcs $J_i$ must lie in the open arc.
\item The case split in the (S4) argument omits a component not inside one arc whose frontier points lie on one arc (around a cycle of $G$); the same chain argument covers it.
\item That the final $S$ irreducibly separates is claimed but not argued, and is not needed: the end sets come out of the construction.
\end{itemize}

\section*{Exercises}

\begin{exercise}
Verify the end sets in Example~\ref{ex:ladder}, and check (S3) and (S4) for them.
\end{exercise}

\begin{exercise}
Construct a ladder on the circle (a triangle in $\R^2$) of height~$3$ with two wedges. Is the result consistent with Claim~14.1?
\end{exercise}

\begin{exercise}
Run the construction of Theorem~\ref{thm:exists-ladder} for $G=\I$, $M=\I\times\set{\frac12}$ and $u(x,t)=t-\frac12$ with $N=4$. What is the ladder?
\end{exercise}

\begin{exercise}
Why are all list lengths odd in Theorem~\ref{thm:exists-ladder}? What goes wrong in the (S4) argument of Theorem~\ref{thm:ladder-stairwell} if lists of even length were allowed?
\end{exercise}
